\section{Struttura degli strumenti}

\subsection{Notion}
Notion viene utilizzato come piattaforma di collaborazione e verifica per i contenuti in fase di produzione. I file vengono redatti in Notion e, dopo la verifica, sono pubblicati sulla repository di documentazione su GitHub.

\subsection{GitHub}
Il progetto MiroFish Trading Offline utilizza due repository:

\begin{enumerate}
  \item \textbf{Repository Codice}: \url{https://github.com/eghosa02/mirofish_trading_offline}
  
  Contiene il codice sorgente del progetto. I branch devono seguire la nomenclatura: \texttt{feature\#<numero\_issue>}
  
  \item \textbf{Repository Documentazione}: \url{https://github.com/eghosa02/mirofish_trading_offline_doc}
  
  Contiene la documentazione ufficiale del progetto.
\end{enumerate}

\subsection{Nomenclatura dei Branch}
Il sistema di versionamento prevede branch con la seguente struttura:
\begin{itemize}
  \item \texttt{feature\#<numero>}: per lo sviluppo di nuove funzionalità
  \item Esempio: \texttt{feature\#42} (per la issue numero 42)
\end{itemize}
